% Teil II: der durchgerechnete Weg.

\subsection*{Gegeben und gesucht}
Gegeben: Durchmesser $d_1,d_2$ der Querschnitte 1 und 2, Höhendifferenz
$\Delta h$ des Quecksilbermanometers, Dichten $\rho$ (Gas) und
$\rho_\mathrm{m}$ (Quecksilber), $g$.
Gesucht: Volumenstrom $\dot V$ durch das Rohr.

\subsection*{Ansatz}
Kontinuität verknüpft die Geschwindigkeiten, Bernoulli die Geschwindigkeits-
mit der Druckdifferenz, das Manometer gibt die Druckdifferenz. Zwei
Unbekannte ($v_1,v_2$), drei Gleichungen --- also nach der einen Geschwindigkeit
auflösen und mit $A_1$ zum Volumenstrom multiplizieren.

\subsection*{Rechnung}
Kontinuität: $v_2 = v_1\frac{A_1}{A_2}$. Einsetzen in Bernoulli
(horizontal, reibungsfrei):
\begin{equation}
  p_1-p_2=\frac{\rho}{2}\,v_1^2\left[\left(\frac{A_1}{A_2}\right)^2-1\right]
  =\frac{\rho}{2}\,v_1^2\left[\left(\frac{d_1}{d_2}\right)^4-1\right]\;.
\end{equation}
Andererseits $p_1-p_2=\rho_\mathrm{m}g\Delta h$, also
\begin{equation}
  v_1=\sqrt{\frac{2\,\rho_\mathrm{m}\,g\,\Delta h}
        {\rho\left[\left(\frac{d_1}{d_2}\right)^4-1\right]}}\;.
\end{equation}
Mit $A_1=\pi d_1^2/4$:

\begin{ergebnis}
\begin{equation*}
  \dot V = A_1 v_1
  =\frac{\pi d_1^2}{4}
   \sqrt{\frac{2\,\rho_\mathrm{m}\,g\,\Delta h}
        {\rho\left[\left(\frac{d_1}{d_2}\right)^4-1\right]}}\;.
\end{equation*}
\end{ergebnis}

\begin{kontrolle}
\textbf{Einheiten:} $\sqrt{\rho_\mathrm{m}g\Delta h/\rho}
\Rightarrow\sqrt{\frac{kg/m^3\cdot m/s^2\cdot m}{kg/m^3}}
=\sqrt{m^2/s^2}=m/s$; mal $m^2$ ergibt $m^3/s$. \checkmark

\textbf{Grenzfall:} Für $d_1=d_2$ (keine Verengung) divergiert der
Volumenstrom --- es gibt keinen Geschwindigkeitsunterschied, das Venturi-Prinzip
versagt. Für $d_2\to 0$ wird $\dot V\to 0$. Plausibel.

\textbf{Größenordnung:} Mit $\Delta h$ im \unit{\centi\metre}-Bereich, Gas
$\rho\!\approx\!\qty{1}{\kg\per\cubic\metre}$, $\rho_\mathrm{m}\!\approx\!\qty{13600}{\kg\per\cubic\metre}$
entsteht ein Volumenstrom im Bereich einiger $\unit{\litre\per\second}$
--- realistisch für ein Venturi-Rohr.
\end{kontrolle}

\begin{merksatz}
Venturi-Messung = Bernoulli (Druck/kinetische Energie) + Kontinuität
(Querschnitt/Geschwindigkeit) + Manometer (Druckdifferenz); die Verengung
übersetzt Strömung in ein messbares Druckzeichen.
\end{merksatz}
